% CODEX: wstaw jako nową sekcję w 30b-geometrizacja-thurstona.tex
% po całym \section{Rozwłóknienia Seiferta i geometria},
% ale przed \section{Twierdzenie geometryzacyjne}.
% Nie powtarzaj definicji \mathbb H^3, modeli Poincarégo ani ogólnych
% twierdzeń o ilorazach. Wyniki głębokie oznaczaj jako zewnętrzne.
\section{Przykłady hiperboliczne: wielościany, ostrogi i wypełnienia}
\label{sec:thur-dop-hiperboliczne}
Model $\mathbb H^3$ opisano już w sekcjach o metryce i o
geometrii Kleina. Teraz skupimy się na tym, jak z tego modelu
zbudować rozmaitości i jak obliczać ich niezmienniki.

\begin{definicja}[Wielościan fundamentalny i sklejanie ścian]
\label{def:thur-dop-wieloscian}
Niech $\Gamma<\operatorname{Isom}(\mathbb H^3)$ będzie grupą
dyskretną. Wielościan $P$ o niepustym wnętrzu nazywamy
\emph{wielościanem fundamentalnym}, gdy wnętrza $\gamma P$
są parami rozłączne dla różnych $\gamma\in\Gamma$,
a ich suma razem z brzegami pokrywa $\mathbb H^3$.
Zadanie odwrotne polega na zadaniu wielościanu i izometrii
parujących jego ściany, a następnie wykazaniu, że wyznaczają
one grupę dyskretną i właściwe sklejanie.
\end{definicja}

Nie wystarczy sam rysunek parowania ścian. W twierdzeniu
Poincarégo o wielościanach trzeba kontrolować między innymi
cykle krawędzi, sumy kątów i zgodność identyfikacji.
Wielościany \emph{idealne} mają wierzchołki na sferze
w nieskończoności $\partial_\infty\mathbb H^3$;
objętość może być skończona, mimo że wielościan
jest niezwarty w metryce hiperbolicznej.

\begin{przyklad}[Czworościany idealne, Gieseking i węzeł ósemkowy]
\label{prz:thur-dop-gieseking-osemka}
Regularny czworościan idealny w $\mathbb H^3$ ma kąty
dwuścienne równe $\pi/3$ i skończoną objętość $v_3$.
Istnieje parowanie jego czterech ścian dające
\emph{rozmaitość Giesekinga}: jest ona nieorientowalna,
zupełna, niezwarta i ma skończoną objętość $v_3$.
Jej jedyna ostroga ma przekrój będący butelką Kleina.
Natomiast odpowiednie sklejenie \emph{dwóch} takich
czworościanów daje dopełnienie węzła ósemkowego
$M_8=S^3\setminus K_8$ o objętości $2v_3$.
Jest ono orientowalne, a jego jedyna ostroga ma przekrój torusowy.
Dopełnienie węzła ósemkowego jest orientowalnym
nakryciem dwukrotnym rozmaitości Giesekinga.

Z punktu widzenia geometrii nie wystarcza sprawdzenie,
że czworościany pokrywają rozmaitość topologiczną:
dookoła każdej klasy krawędzi suma kątów musi wynosić
$2\pi$, a identyfikacje wierzchołków idealnych
muszą dać zupełne przekroje euklidesowe.
Dla regularnego czworościanu sześć kątów po $\pi/3$
daje właśnie $2\pi$; w konkretnych sklejaniach należy
jeszcze sprawdzić wszystkie klasy krawędzi i wierzchołków.
Istnienie opisanych sklejeń i ich własności traktujemy
tutaj jako znane wyniki, nie jako konsekwencję samej sumy kątów.
\end{przyklad}


\begin{przyklad}[Zamknięta przestrzeń dwunastościenna Seiferta--Webera]
\label{prz:thur-dop-seifert-weber}
W poprzednim przykładzie otrzymaliśmy rozmaitości niezwarte
z wierzchołków idealnych. Istnieją także zwarte przykłady
skonstruowane z pojedynczego \emph{zwartego} wielościanu.
Weźmy foremny dwunastościan hiperboliczny o kącie dwuściennym
$2\pi/5$. Sklejamy każdą parę przeciwległych pięciokątnych
ścian, wykonując obrót o $3/10$ pełnego obrotu, czyli $3\pi/5$.
Pod tymi identyfikacjami trzydzieści krawędzi tworzy
sześć klas po pięć krawędzi. Suma kątów wokół każdej
krawędzi ilorazu wynosi więc $5(2\pi/5)=2\pi$.
Sklejenie jest zgodne również w wierzchołkach; jego
rezultatem jest zamknięta orientowalna rozmaitość
hiperboliczna, zwana przestrzenią Seiferta--Webera.
Kontrola otoczeń wierzchołków i pełnych warunków
sklejenia to istotna część konstrukcji, którą
przywołujemy tutaj jako znany wynik.
Nie należy mylić jej ze sferyczną przestrzenią
dwunastościenną Poincarégo: powstaje ona z innego
wielościanu i innej identyfikacji ścian.
\end{przyklad}

\begin{twierdzenie}[Sztywność Mostowa--Prasada]
\label{thm:thur-dop-mostow-prasad}
Niech $M$ i $N$ będą zupełnymi rozmaitościami
hiperbolicznymi o skończonej objętości i tym samym
wymiarze $n\ge3$. Każdy izomorfizm
$\pi_1(M)\cong\pi_1(N)$ jest realizowany przez izometrię.
W szczególności objętość hiperboliczna jest
niezmiennikiem topologicznym takich rozmaitości.
\end{twierdzenie}
Wynik ten cytujemy bez dowodu. Jego sens różni się od
sytuacji powierzchni hiperbolicznych, gdzie ustalona
topologia zazwyczaj dopuszcza nietrywialne deformacje
metryk o krzywiźnie stałej $-1$.

\begin{definicja}[Wypełnienie Dehna]
\label{def:thur-dop-dehn-fill}
Niech $N$ będzie zwartą orientowalną trójrozmaitością
z brzegiem zawierającym torus $T$. Wybierzmy
uporządkowaną bazę $(\mu,\lambda)$ grupy
$H_1(T;\Z)\cong\Z^2$. \emph{Nachylenie} jest klasą
niezorientowanej prostej zamkniętej na $T$ o prymitywnej
klasie homologii $p\mu+q\lambda$, $\gcd(p,q)=1$.
Wypełnienie Dehna $N(p/q)$ otrzymujemy przez
przyklejenie pełnego torusa $D^2\times S^1$ do $T$
w taki sposób, aby obraz brzegu dysku meridianowego
$\partial D^2\times\{*\}$ przedstawiał wybrane nachylenie.
\end{definicja}

\begin{stwierdzenie}[Grupa podstawowa po wypełnieniu]
\label{prop:thur-dop-dehn-pi1}
Niech $N$ ma jedną składową brzegu torusowego,
a $\mu,\lambda$ oznaczają też elementy $\pi_1(N)$
wyznaczone przez ich obrazy w $N$, z ustalonymi
drogami bazowymi. Wówczas
\[
 \pi_1(N(p/q))\cong
 \pi_1(N)/\langle\!\langle\mu^p\lambda^q\rangle\!\rangle.
\]
\end{stwierdzenie}
\begin{proof}
Do $N$ dołączamy pełny torus, którego meridian staje się
ściągalny. Twierdzenie van Kampena identyfikuje grupę
przestrzeni sklejonej z ilorazem $\pi_1(N)$ przez
normalne domknięcie tej krzywej meridianowej.
Ponieważ krzywe $\mu,\lambda$ na torusie komutują,
klasa nachylenia ma w grupie postać $\mu^p\lambda^q$.
\end{proof}

\begin{przyklad}[Homologia wypełnienia dopełnienia węzła]
\label{prz:thur-dop-dehn-h1}
Dla dopełnienia węzła $K\subset S^3$ mamy
$H_1(S^3\setminus\operatorname{int}\nu K;\Z)\cong\Z$,
generowaną przez klasę meridianu $\mu$; preferowana
longitude $\lambda$ znika w tej grupie. Po abelianizacji
relacji z poprzedniego stwierdzenia dostajemy
\[
 H_1(N(p/q);\Z)\cong
 \begin{cases}\Z/|p|\Z,&p\ne0,\\\Z,&p=0.\end{cases}
\]
Przykład pokazuje, jak operacja geometryczna zmienia
niezmiennik topologiczny. Nie wynika z niego jednak,
że dla $p=\pm1$ otrzymana rozmaitość jest sferą $S^3$:
zerowa pierwsza homologia nie rozstrzyga grupy podstawowej.
\end{przyklad}

\begin{twierdzenie}[Thurston: hiperboliczne wypełnienia Dehna]
\label{thm:thur-dop-dehn-hyperbolic}
Jeżeli wnętrze zwartej orientowalnej trójrozmaitości $N$
z jedną torusową składową brzegu ma zupełną metrykę
hiperboliczną o skończonej objętości, to wszystkie poza
skończenie wieloma nachyleniami dają zamknięte
trójrozmaitości hiperboliczne $N(p/q)$.
\end{twierdzenie}
Twierdzenie to jest wynikiem zewnętrznym Thurstona,
a nie następstwem samej prezentacji grupy podstawowej.
Pokazuje ono, że z jednej niezwartej rozmaitości
hiperbolicznej otrzymujemy nieskończenie wiele
zamkniętych przykładów.

\begin{uwaga}[Związek z geometrią Sol]
Dla węzła ósemkowego nachylenie zerowe daje
szczególne, niehiperboliczne wypełnienie: torus odwzorowania
hiperbolicznego automorfizmu $T^2$, a więc rozmaitość
geometrii Sol. To łączy geometrię węzłów z przykładem
monodromii macierzy hiperbolicznej w Sekcji~\ref{sec:thurston-eight}.
Sformułowanie o skończenie wielu wyjątkach nie pozwala jednak
samodzielnie ustalić listy wszystkich wyjątkowych nachyleń.
\end{uwaga}

\par\smallskip
{\footnotesize\noindent\emph{Dalsza lektura:}
\href{https://www.math.uic.edu/~kauffman/JessicaPurcell.pdf}
{J.~S.~Purcell, \emph{Hyperbolic Knot Theory}};
\href{https://people.dm.unipi.it/martelli/Geometric_topology.pdf}
{B.~Martelli, \emph{An Introduction to Geometric Topology}},
rozdziały o geometriach Sol i hiperbolicznych rozmaitościach
z ostrogami; S.~A.~Pinillos Prado, \emph{Geometric Structures on
3-Dimensional Manifolds} (2021), rozdziały 2--3.\par}
